%%
%% part3-customization/ch22-localization.tex
%%
%% Chapter 22: Localization — Persian (fa-IR) and English
%% (en-US) locales, language switching, and number/date
%% formatting.
%%

\chapter{لوکال و زبان}
\label{chap:localization}

\section{چرا این موضوع مهم است؟}
\label{sec:localization-why}

\lr{MatinBook} برای کتاب‌های فنی فارسی طراحی
شده است، ولی از کتاب‌های انگلیسی نیز پشتیبانی
می‌کند. لوکال‌ها (فایل‌های ترجمه) نام عناصر
مختلف کتاب را به زبان مورد نظر ترجمه می‌کنند.

در این فصل، با دو لوکال \lr{fa-IR} و
\lr{en-US} آشنا می‌شوید و یاد می‌گیرید که
چگونه بین آن‌ها سوئیچ کنید.

\mnote{لوکال \lr{fa-IR} به‌طور پیش‌فرض در \lr{MatinBook} بارگذاری می‌شود. لوکال \lr{en-US} برای کتاب‌های انگلیسی است.}

\section{لوکال فارسی (\lr{fa-IR})}
\label{sec:localization-fa}

لوکال \lr{fa-IR} نام عناصر \lr{LaTeX} را
به فارسی ترجمه می‌کند. جدول
\cref{tab:fa-names} این ترجمه‌ها را نشان
می‌دهد.

\begin{matintable}{نام‌های فارسی عناصر \lr{LaTeX}}{tab:fa-names}
\begin{tblr}{
    width = \textwidth,
    colspec = {l X[l] X[l]},
    hlines, vlines,
    colsep = 6pt,
    rowsep = 4pt,
    row{1} = {font=\bfseries},
}
نام \lr{LaTeX} & ترجمه‌ی فارسی & کاربرد \\
\lr{contentsname} & فهرست مطالب & فهرست \\
\lr{listfigurename} & فهرست تصاویر & فهرست شکل‌ها \\
\lr{listtablename} & فهرست جداول & فهرست جدول‌ها \\
\lr{bibname} & منابع و مراجع & کتاب‌نامه \\
\lr{indexname} & نمایه & نمایه \\
\lr{chaptername} & فصل & سرصفحه \\
\lr{appendixname} & پیوست & پیوست‌ها \\
\lr{proofname} & اثبات & محیط اثبات \\
\lr{figurename} & شکل & کپشن شکل \\
\lr{tablename} & جدول & کپشن جدول \\
\end{tblr}
\end{matintable}

\mnote{این نام‌ها در فایل \lr{fa-IR.sty} تعریف شده‌اند و با \lr{\textbackslash renewcommand} بازنویسی می‌شوند.}

\subsection{نام‌های \lr{cleveref}}
\label{sec:localization-fa-cref}

لوکال فارسی، نام‌های \lr{cleveref} را نیز
ترجمه می‌کند. جدول \cref{tab:fa-cref-names}
این ترجمه‌ها را نشان می‌دهد.

\begin{matintable}{نام‌های فارسی \lr{cleveref}}{tab:fa-cref-names}
\begin{tblr}{
    width = \textwidth,
    colspec = {l X[l] X[l]},
    hlines, vlines,
    colsep = 6pt,
    rowsep = 4pt,
    row{1} = {font=\bfseries},
}
نوع & مفرد & جمع \\
معادله & معادله & معادلات \\
فصل & فصل & فصول \\
بخش & بخش & بخش‌ها \\
شکل & شکل & اشکال \\
جدول & جدول & جداول \\
الگوریتم & الگوریتم & الگوریتم‌ها \\
قضیه & قضیه & قضایا \\
لم & لم & لم‌ها \\
نتیجه & نتیجه & نتایج \\
تعریف & تعریف & تعاریف \\
مثال & مثال & مثال‌ها \\
نکته & نکته & نکات \\
تمرین & تمرین & تمرین‌ها \\
گزاره & گزاره & گزاره‌ها \\
\end{tblr}
\end{matintable}

\mnote{این نام‌ها با \lr{\textbackslash crefname} تعریف می‌شوند و در \lr{mb-theorem.sty} و \lr{fa-IR.sty} تنظیم شده‌اند.}

\subsection{نام‌های محیط‌های علمی}
\label{sec:localization-fa-theorems}

نام‌های محیط‌های علمی نیز در لوکال فارسی
تعریف شده‌اند:

\begin{itemize}
    \item \cmd{theoremname}: قضیه.
    \item \cmd{lemmaname}: لم.
    \item \cmd{corollaryname}: نتیجه.
    \item \cmd{propositionname}: گزاره.
    \item \cmd{definitionname}: تعریف.
    \item \cmd{examplename}: مثال.
    \item \cmd{remarkname}: نکته.
    \item \cmd{exercisename}: تمرین.
    \item \cmd{solutionname}: راه‌حل.
\end{itemize}

\mnote{این نام‌ها در عنوان جعبه‌های علمی استفاده می‌شوند و در فایل \lr{fa-IR.sty} تعریف شده‌اند.}

\section{لوکال انگلیسی (\lr{en-US})}
\label{sec:localization-en}

لوکال \lr{en-US} نسخه‌ی انگلیسی
\lr{fa-IR} است و برای کتاب‌های انگلیسی
استفاده می‌شود.

\begin{latin}
\begin{minted}{latex}
% |\pc{در فایل}|\lr{matinbook.cls}:
% |\pc{به‌جای}|%
\RequirePackage{fa-IR}

% |\pc{بنویسید}|%
\RequirePackage{en-US}
\end{minted}
\end{latin}

\mnote{لوکال \lr{en-US} از ساختار مشابه \lr{fa-IR} پیروی می‌کند ولی ترجمه‌ها انگلیسی هستند.}

\section{سوئیچ بین زبان‌ها}
\label{sec:localization-switching}

\lr{MatinBook} در حال حاضر از سوئیچ زبان
در زمان اجرا پشتیبانی نمی‌کند. یعنی
نمی‌توانید در وسط سند، زبان را از فارسی به
انگلیسی تغییر دهید.

برای سوئیچ، باید کلاس را با لوکال انگلیسی
کامپایل کنید:

\begin{latin}
\begin{minted}{latex}
% |\pc{در فایل}|\lr{matinbook.cls}:
\RequirePackage{en-US}
\end{minted}
\end{latin}

\mnote{پشتیبانی از سوئیچ زبان در زمان اجرا، یکی از برنامه‌های آینده‌ی \lr{MatinBook} است.}

\section{متن دوزبانه}
\label{sec:localization-bilingual}

در کتاب‌های دوزبانه، می‌توانید متن فارسی و
انگلیسی را با هم ترکیب کنید:

\begin{latin}
\begin{minted}{latex}
|\pc{این کتاب درباره}|\lr{machine learning}|\pc{و}|\lr{deep learning}|\pc{است.}|%
\end{minted}
\end{latin}

خروجی: این کتاب درباره \lr{machine learning}
و \lr{deep learning} است.

\mnote{برای متن لاتین در متن فارسی، از \lr{\textbackslash lr} استفاده کنید. برای پاراگراف کامل لاتین، از محیط \lr{latin}.}

\section{قالب اعداد}
\label{sec:localization-numbers}

\subsection{اعداد در متن فارسی}
\label{sec:localization-numbers-fa}

در متن فارسی، اعداد به‌طور خودکار به ارقام
فارسی (۰-۹) تبدیل می‌شوند:

\begin{latin}
\begin{minted}{latex}
|\pc{این کتاب شامل ۱۲ فصل است.}|%
\end{minted}
\end{latin}

خروجی: این کتاب شامل ۱۲ فصل است.

\subsection{اعداد در متن لاتین}
\label{sec:localization-numbers-latin}

در متن لاتین، اعداد به‌صورت لاتین (0-9)
نمایش داده می‌شوند:

\begin{latin}
\begin{minted}{latex}
\begin{latin}
The book has 12 chapters.
\end{latin}
\end{minted}
\end{latin}

خروجی:

\begin{latin}
The book has 12 chapters.
\end{latin}

\subsection{اعداد اعشاری فارسی}
\label{sec:localization-numbers-decimal}

برای اعداد اعشاری فارسی، از دستور \cmd{pnum}
استفاده کنید:

\begin{latin}
\begin{minted}{latex}
|\pc{مقیاس}|\pnum{۱.۲}|\pc{به این معناست...}|%
\end{minted}
\end{latin}

خروجی: مقیاس \pnum{۱.۲} به این معناست...

\mnote{دستور \lr{\textbackslash pnum} از باگ معکوس‌شدن اعداد اعشاری در \lr{bidi} جلوگیری می‌کند.}

\subsection{اعداد در ریاضیات}
\label{sec:localization-numbers-math}

در ریاضیات، اعداد از بافت پیروی می‌کنند:

\begin{itemize}
    \item در متن فارسی: اعداد ریاضی به‌صورت
    فارسی نمایش داده می‌شوند.
    \item در \cmd{lr}: اعداد ریاضی به‌صورت
    لاتین نمایش داده می‌شوند.
\end{itemize}

\begin{latin}
\begin{minted}{latex}
|\pc{معادله}|\lr{$E = mc^{2}$}|\pc{...}|%
\end{minted}
\end{latin}

\mnote{برای نمایش اعداد لاتین در ریاضیات، از \lr{\textbackslash lr\{...\}} استفاده کنید.}

\section{قالب تاریخ}
\label{sec:localization-dates}

\subsection{تاریخ فارسی}
\label{sec:localization-dates-fa}

برای تاریخ فارسی، از تاریخ تقویم هجری
شمسی استفاده کنید:

\begin{latin}
\begin{minted}{latex}
\date{|\pc{مهر ۱۴۰۵}|}
\end{minted}
\end{latin}

خروجی: مهر ۱۴۰۵.

\mnote{دستور \lr{\textbackslash today} در حالت فارسی، تاریخ هجری شمسی را نمایش می‌دهد.}

\subsection{تاریخ انگلیسی}
\label{sec:localization-dates-en}

در حالت انگلیسی، از تاریخ میلادی استفاده
کنید:

\begin{latin}
\begin{minted}{latex}
\date{October 2026}
\end{minted}
\end{latin}

\section{شخصی‌سازی لوکال}
\label{sec:localization-custom}

برای شخصی‌سازی نام‌ها، فایل \file{fa-IR.sty}
یا \file{en-US.sty} را ویرایش کنید. مثلاً
برای تغییر «فهرست مطالب» به «محتویات»:

\begin{latin}
\begin{minted}{latex}
% |\pc{در فایل}|\lr{fa-IR.sty}:
\renewcommand{\contentsname}{|\pc{محتویات}|}
\end{minted}
\end{latin}

\mnote{برای تغییرات موقت، می‌توانید در سند خود نیز این دستورات را بازنویسی کنید.}

\section{خطاهای رایج}
\label{sec:localization-mistakes}

\subsection{فراموش کردن \cmd{lr}}
\label{sec:localization-mistake-lr}

\textbf{مشکل:} اگر متن لاتین را در متن فارسی
بدون \cmd{lr} بنویسید، جهت متن به‌هم
می‌ریزد.

\textbf{راه‌حل:} همیشه از \cmd{lr} برای متن
لاتین در متن فارسی استفاده کنید.

\subsection{فراموش کردن \cmd{pnum}}
\label{sec:localization-mistake-pnum}

\textbf{مشکل:} اگر اعداد اعشاری فارسی را
بدون \cmd{pnum} بنویسید، به‌صورت معکوس
نمایش داده می‌شوند (مثلاً \lr{۱.۲} می‌شود
\lr{۲.۱}).

\textbf{راه‌حل:} از \cmd{pnum} برای همه‌ی
اعداد اعشاری فارسی استفاده کنید.

\subsection{ترکیب لوکال‌ها}
\label{sec:localization-mistake-mixed}

\textbf{مشکل:} اگر هر دو لوکال \lr{fa-IR} و
\lr{en-US} را بارگذاری کنید، لوکال دوم
اولی را بازنویسی می‌کند.

\textbf{راه‌حل:} فقط یکی از لوکال‌ها را
بارگذاری کنید.

\section{تمرین‌ها}
\label{sec:localization-exercises}

\begin{exercise}
    یک کتاب با لوکال \lr{en-US} بسازید و
    ببینید که نام‌ها به انگلیسی نمایش
    داده می‌شوند.
    \label{exc:localization-en}
\end{exercise}

\begin{exercise}
    نام «فهرست مطالب» را به «محتویات» تغییر
    دهید.
    \label{exc:localization-custom}
\end{exercise}

\begin{exercise}
    یک پاراگراف فارسی با سه واژه‌ی لاتین
    بنویسید و از \cmd{lr} استفاده کنید.
    \label{exc:localization-lr}
\end{exercise}

\begin{exercise}
    یک عدد اعشاری فارسی را با و بدون
    \cmd{pnum} بنویسید و تفاوت را ببینید.
    \label{exc:localization-pnum}
\end{exercise}

\section{خلاصه}
\label{sec:localization-summary}

در این فصل، با لوکال و زبان در \lr{MatinBook}
آشنا شدیم:

\begin{itemize}
    \item \lr{MatinBook} از دو لوکال پشتیبانی
    می‌کند: \lr{fa-IR} (فارسی) و \lr{en-US}
    (انگلیسی).

    \item لوکال فارسی، نام‌های \lr{LaTeX} و
    \lr{cleveref} را به فارسی ترجمه می‌کند.

    \item لوکال انگلیسی، نسخه‌ی انگلیسی
    \lr{fa-IR} است.

    \item سوئیچ بین زبان‌ها با تغییر
    \cmd{RequirePackage\{fa-IR\}} به
    \cmd{RequirePackage\{en-US\}} در
    \file{matinbook.cls} انجام می‌شود.

    \item در متن دوزبانه، از \cmd{lr} برای
    متن لاتین در متن فارسی استفاده کنید.

    \item اعداد در متن فارسی به‌طور خودکار
    فارسی می‌شوند.

    \item برای اعداد اعشاری فارسی، از
    \cmd{pnum} استفاده کنید.

    \item تاریخ فارسی با تقویم هجری شمسی.

    \item خطاهای رایج شامل فراموش کردن
    \cmd{lr}، فراموش کردن \cmd{pnum}، و
    ترکیب لوکال‌ها است.
\end{itemize}

این فصل، آخرین فصل از بخش سوم (شخصی‌سازی)
بود. در بخش چهارم (مباحث پیشرفته)، با
ساختار ماژولار، تم سفارشی، دستورات سفارشی
و پروژه‌های بزرگ آشنا می‌شوید.